% Unrefereed preprint for public release.
% Copyright 2026 Ricardo Pinho. CC BY 4.0.
\documentclass[11pt]{amsart}
\usepackage[T1]{fontenc}
\usepackage{lmodern}
\usepackage[margin=1in]{geometry}
\usepackage{amsmath,amssymb,booktabs}
\usepackage[expansion=false]{microtype}
\setlength{\emergencystretch}{2em}
\usepackage[hyphens]{url}
\usepackage[hidelinks]{hyperref}
\newtheorem{theorem}{Theorem}
\newtheorem{lemma}[theorem]{Lemma}
\newcommand{\D}{\mathcal D}
\title[A local-degree cap]{A degree-three local-weight cap and an improved relevant-variable bound}
\author{Ricardo Pinho}
\address{Independent researcher}
\thanks{ORCID: \href{https://orcid.org/0009-0006-2744-2599}{0009-0006-2744-2599}. Licensed under CC BY 4.0.}
\date{23 October 2026; unrefereed preprint}
\keywords{Boolean functions, relevant variables, polynomial degree, exact arithmetic}
\begin{document}
\begin{abstract}
We prove that the local-degree potential of a Boolean function of degree
at most three is at most $5/4$, using the published fact that such a
function has at most ten relevant variables. Inserting this cap into
Wellens's degree--block-sensitivity recurrence, with the existing exact
certificates released by \texorpdfstring{\'Epi}{Epi}, gives
$C_{44}\le4.3576549226$.
The reproduced preceding recurrence gives approximately $4.365467422574458$;
our exact degree-three cap lowers it by $1/128$.
Two separately written exact-arithmetic programs check the finite
certificates, recurrence, infinite tail and finite cases of the added lemma.
The universal argument is written below; the published degree-three
classification is cited, not re-enumerated. This preprint is not peer
reviewed and not verified in a proof assistant; no optimality is claimed.
\end{abstract}
\maketitle

\section{The constant and the preceding bound}
For a Boolean function $f$, let $r(f)$ be its number of relevant variables
and $d=\deg f$ the degree of its multilinear polynomial. We use the sign
cube and sign-valued output when convenient; invertible affine changes of
input and output preserve degree and local degrees. The constant is
\[
 C_{44}=\sup_{d\ge1}\ \sup_{\deg f\le d}\frac{r(f)}{2^d}.
\]
Write $d_i(f)$ for the largest degree of a monomial involving a relevant
coordinate $i$, and
\[
 \D(f)=\sum_{i\text{ relevant}}2^{-d_i(f)}.
\]
Since $d_i(f)\le d$, $r(f)/2^d\le\D(f)$.

\begin{theorem}\label{main}
For every Boolean function $f$ of positive degree $d$,
\[
 r(f)\le \frac{21788274613}{5000000000}\,2^d
       =4.3576549226\,2^d.
\]
Consequently $C_{44}\le4.3576549226$.
\end{theorem}

The merged table records $4.394$ from Wellens~\cite{wellens,table}.
The stronger public bound located in the present audit is \'Epi's
$4.3654674227$~\cite{epi,epipaper}. Its degree-48 certificates and
horizon-96 recurrence reproduce exactly at
\[
 \frac{172933981203227910490173314851}
      {39614081257132168796771975168}
 =4.365467422574458\ldots.
\]
That extra decimal precision is not a new improvement. We keep these
released polynomial certificates and change one analytic cap in the
recurrence. The difference between the two exact horizon-96 bounds is
$1/128$, a gain in the actual preceding method rather than a rounding gain.

\section{A sharp degree-three potential cap}
We recall two ingredients. First, the influence argument gives
\begin{equation}\label{influence}
 \D(f)\le \tfrac12 I(f)\le \tfrac12\deg f.
\end{equation}
Indeed, a nonzero derivative of local degree $d_i-1$ is nonzero on at least
a $2^{-(d_i-1)}$ fraction of the cube; summing gives the first inequality.
The Fourier expression for total influence gives the second.

Second, for every set $H$ of coordinates, the restriction inequality is
\begin{equation}\label{restrict}
 \D(f)\le\sum_{i\in H\text{ relevant}}2^{-d_i(f)}
             +\mathbb E_{\alpha\in\{-1,1\}^H}\D(f_\alpha).
\end{equation}
For one restricted coordinate, any other coordinate surviving in both
branches has local degree no greater than its original local degree. If
it survives in only one branch, its original local degree is exactly one
higher than in that branch. Thus its original weight is at most the
average restricted weight in either case. Sum over the unrestricted
coordinates and iterate to get~\eqref{restrict}. These are the usual
local-weight facts of Chiarelli--Hatami--Saks and Wellens~\cite{chs,wellens}.

We use the established theorem $R_3=10$: a Boolean function of degree at
most three depends on at most ten variables. This comes from
Tarannikov--Kirienko's resilient-function result and its degree translation;
Krotov--Valyuzhenich explicitly state this consequence in Section 3 and
step (i) of their classification~\cite{tk,kv}. Their account also cites
Zverev's computer-free argument. We rely on this published theorem;
our programs do not independently re-enumerate all cubic Boolean functions.

\begin{lemma}\label{cap}
If $\deg f\le3$, then $\D(f)\le5/4$. The constant $5/4$ is attained.
\end{lemma}
\begin{proof}
Degrees at most two satisfy $\D(f)\le1$ by~\eqref{influence}, so assume
$\deg f=3$. A relevant coordinate of local degree one would give
$f=x_i a+g$, with $a$ a nonzero constant. Both $g+a$ and $g-a$ must take
values in $\{-1,1\}$, forcing $a=\pm1$ and $g=0$. Such a function is a
dictator, contrary to degree three. Hence all relevant local degrees are
at least two.

If they are all three, the published ten-variable theorem immediately gives
$\D(f)=r(f)/8\le10/8=5/4$. Otherwise fix a coordinate $i$ of local degree
two. Write
\[
 f(x)=x_i A(y)+B(y),\qquad
 A(y)=\frac{f(1,y)-f(-1,y)}2.
\]
The nonzero polynomial $A$ is affine and takes values only in $\{-1,0,1\}$.
Every nonzero linear coefficient $b_j$ has $|b_j|\ge1/2$, since changing
that input changes $A$ by $2b_j$, a nonzero integer. Its total range has
width $2\sum_j|b_j|\le2$. Thus $A$ depends on a set $H$ of at most two
coordinates. Let $s=|H|$ and $p=\Pr(A\ne0)$ over those controls.

Whenever $A=\pm1$, Booleanity forces $B=0$, so the corresponding restricted
function is $\pm x_i$, with potential $1/2$. Whenever $A=0$, the restricted
function is independent of $i$ and has degree at most two. To see the degree
claim, a cubic monomial disjoint from $H$ would have the same coefficient
in every restriction to $H$: a term containing that monomial and an additional
control would have degree at least four. But one of the nonempty active
branches is a dictator, which forces each such cubic coefficient to vanish.
All restrictions therefore have degree at most two. Their potentials are
at most one by~\eqref{influence}.

Each coordinate in $H$ has original local degree at least two. Applying
\eqref{restrict} gives
\[
 \D(f)\le\frac{s}{4}+\frac p2+(1-p).
\]
The possible nonzero affine ternary forms are, up to signs, coordinate
permutations and omission of unused controls: $1$, $x_j$,
$(1+x_j)/2$, and $(x_j+x_k)/2$.
For $(s,p)=(0,1),(1,1),(1,1/2),(2,1/2)$, respectively, the last bound is
$1/2,3/4,1,5/4$. This proves the inequality.

For sharpness define
\[
 Q(s,t,u,v)=\tfrac12\bigl((s+t)u+(s-t)v\bigr),
\]
and, on disjoint blocks $x,y$ of four variables each, set
\[
 F(s,t,x,y)=\tfrac12\bigl((s+t)Q(x)+(s-t)Q(y)\bigr).
\]
For sign inputs exactly one summand is active in each expression, so $F$
is Boolean. Its sixteen nonzero monomials are cubic, and each of its ten
variables appears in one. Hence $\D(F)=10/8=5/4$. Both checkers independently
verify this small witness.
\end{proof}

\section{The exact finite recurrence and tail}
Let $\mathcal B_d$ be an upper bound for the block sensitivity of degree-$d$
Boolean functions. Wellens's restriction of a degree-$d$ monomial costs
$d2^{-d}$ in potential and lowers block sensitivity by at least one.
With $U(0,d)=0$, a valid recursion is
\begin{equation}\label{dp}
 U(b,d)=\min\!\left\{c_d,\ d2^{-d}+
                  \max_{1\le k\le d}U(b-1,k)\right\}
 \quad(1\le b\le\mathcal B_d),
\end{equation}
where $U(b,d)=U(\mathcal B_d,d)$ for $b>\mathcal B_d$ and
\[
 c_d=d/2\quad(d\ne3),\qquad c_3=5/4.
\]
Induction on $b$ proves that $U(b,d)$ bounds the potential of every
function with degree exactly $d$ and block sensitivity at most $b$.
Constants contribute zero. This is Wellens's recurrence with
Lemma~\ref{cap} replacing the degree-three influence cap $3/2$.

For $d=1,\ldots,48$ we retain \'Epi's released nodes
$1=x_0<x_1<\cdots<x_d$ and rational weights $w_i$.
The exact conditions checked are
\[
 \sum_{i=0}^d w_i x_i^j=0\ (j=1,\ldots,d),\qquad
 w_0>0,\qquad w_0+\sum_{i\ge1:w_i<0}w_i>0.
\]
If a degree-$d$ Boolean function had block sensitivity at least $x_d$,
Wellens's symmetrization criterion would give a polynomial $p$ with
$p(0)=0$, $p(1)=1$ and $0\le p(k)\le1$ for $1\le k\le x_d$.
The displayed identities would imply both $\sum_iw_i p(x_i)=0$ and
$\sum_iw_i p(x_i)>0$, a contradiction. Thus
$\mathcal B_d=x_d-1$ is valid. We intersect these bounds with Wellens's
universal inequality
\[
 3b(b-1)\le2(d^4-d^2),
\]
using the largest integer satisfying it at each higher degree as well.
No floating linear-program status enters this certificate.

Wellens's Corollary 3.10, equations (3.14)--(3.15), supplies the infinite
extension: for every horizon $H$,
\begin{equation}\label{tail}
 \D(f)\le U(\mathcal B_H,H)
       +\frac{(H+1)^3}{2^{H+1}}
       +\sum_{k=H+2}^\infty\frac{2k^2-k}{2^k}.
\end{equation}
The tail in~\eqref{tail} is exactly
\[
 T_H=\frac{(H+1)^3+2H^2+11H+19}{2^{H+1}}.
\]
Our row-wise dyadic-integer computation at $H=96$ gives
\[
 U(\mathcal B_{96},96)+T_{96}
 =\frac{172624496193406565421448533795}
       {39614081257132168796771975168}
 <\frac{21788274613}{5000000000}.
\]
Together with $r(f)2^{-d}\le\D(f)$, this proves Theorem~\ref{main}.

\section{Reproducible verification and scope}
The original \'Epi artifact \texttt{sources/wellens-certificate.json} has
SHA-256
\begin{center}\small
\texttt{ce66f1097a7284e9845113fa1dc39267932cc934fd23d9019ff13d3dcb39cf49}.
\end{center}
The new small input is \texttt{certificates/c44\_candidate.json}.
The first checker tests all power-moment identities and uses a row-wise
integer recurrence at horizon 96. The second verifies the different
Lagrange-residue identities
$w_i x_i\prod_{j\ne i}(x_i-x_j)=\text{constant}$, then computes by columns
using rational fractions and the independent horizon 72. It sums the tail
from separate geometric-series moment formulas. Both prove the same
printed upper bound. Their finite derivative-case tests enumerate
coefficient patterns and affine ternary truth tables, respectively;
the ten-variable sharpness witness is checked by an integer Walsh transform
and by direct evaluation of its cubic polynomial.

Run
\begin{verbatim}
python verify_c44_exact.py certificates/c44_candidate.json
python independent/verify_c44_fraction.py \
       certificates/c44_candidate.json
python rejection_controls.py
\end{verbatim}
Controls reject a smaller unsupported cap, an excessive upper-bound
claim, an altered external theorem, changed source digest, and corrupted
weights or nodes even when the attacker also refreshes the digest.
These programs were written separately, with no imports from search or
from one another, by the same AI agent. They are algorithmic cross-checks;
there has been no independent human review. They do not formalize the
published ten-variable theorem, Wellens's universal recurrence argument,
or the full prose proof of Lemma~\ref{cap}.

Endpoint-refined LP searches through degree 40 and integer slice-count
relaxations through degree 10 did not yield a stronger block-sensitivity
cap. Those are bounded negative experiments, not proofs of optimality.
Their failure prompted the local-weight lemma; the existing polynomial
certificates themselves were not improved.

\section*{Provenance and limitations}
The bounded live-source audit and final refresh are in
\texttt{PRIOR\_ART.md} and \texttt{FINAL\_LIVE\_CHECK.md}.
We build on the local potential of Chiarelli--Hatami--Saks, the recurrence,
symmetrization and tail of Wellens, \'Epi's exact degree certificates, and
the published ten-variable theorem. The contribution is the sharpened
weighted degree-three cap and its effect on that recurrence.
No stronger C44 bound was located in the sources checked; exhaustive
priority is not asserted. This AI-prepared preprint is not peer reviewed
and not verified in a proof assistant. No optimality is claimed.

Analysis and proof drafting by AI agents under my direction; I reviewed
the results at summary level. AI agents also prepared the certificate,
checkers and publication materials. This does not claim a line-by-line
human verification.

\begin{thebibliography}{9}
\bibitem{table} T. Tao and contributors, \emph{Optimization Constants in
Mathematics}, entry 44a,
\url{https://teorth.github.io/optimizationproblems/constants/44a.html}.
\bibitem{chs} J. Chiarelli, P. Hatami and M. Saks,
An asymptotically tight bound on the number of relevant variables in a
bounded degree Boolean function, \emph{Combinatorica} \textbf{40} (2020),
237--244; \url{https://arxiv.org/abs/1801.08564}.
\bibitem{wellens} J. Wellens, Relationships between the number of inputs
and other complexity measures of Boolean functions,
\emph{Discrete Analysis} 2022:19;
\url{https://arxiv.org/html/2005.00566v2}.
\bibitem{epi} BakeLab, \emph{\'Epi-results: Wellens' Boolean-function bound},
public exact certificate, reviewed 11 September 2026,
\url{https://github.com/BakeLab/Epi-results/tree/97fa6a420f417458e11723c8bceb70f1529e447a/mathematics/wellens}.
\bibitem{epipaper} Y. Huang et al., \emph{\'Epi: Towards a Persistent,
Self-Improving Agent for Autonomous Research}, 2026, Appendix A.1,
\url{https://arxiv.org/abs/2610.24944}.
\bibitem{tk} Yu. Tarannikov and D. Kirienko, \emph{Spectral analysis of
high order correlation immune functions}, IACR ePrint 2000/050,
\url{https://eprint.iacr.org/2000/050}.
\bibitem{kv} D. S. Krotov and A. Valyuzhenich,
On degree-3 and $(n-4)$-correlation-immune perfect colorings of $n$-cubes,
\emph{Discrete Mathematics} \textbf{347} (2024), 114138;
\url{https://arxiv.org/html/2311.05566v2}.
\end{thebibliography}
\end{document}
